\documentclass[12pt,a4paper]{article}

\usepackage[T1]{fontenc}
\usepackage[utf8]{inputenc}
\usepackage[brazil]{babel}
\usepackage{lmodern}
\usepackage{geometry}
\usepackage{setspace}
\usepackage{hyperref}
\usepackage{booktabs}
\usepackage{longtable}
\usepackage{listings}
\usepackage{xcolor}

\geometry{margin=2.5cm}
\onehalfspacing

\hypersetup{
    colorlinks=true,
    linkcolor=black,
    urlcolor=blue,
    citecolor=black
}

\lstdefinestyle{cppstyle}{
    language=C++,
    basicstyle=\ttfamily\small,
    keywordstyle=\color{blue},
    commentstyle=\color{gray},
    stringstyle=\color{red!60!black},
    numbers=left,
    numberstyle=\tiny\color{gray},
    stepnumber=1,
    numbersep=8pt,
    showstringspaces=false,
    breaklines=true,
    frame=single
}

\title{Relatorio Tecnico\\Implementacao de Grafo Direcionado com Leitura por Arquivo}
\author{NOME DA ALUNA \\ Curso de Ciencia da Computacao \\ PUC}
\date{\today}

\begin{document}

\maketitle
\tableofcontents
\newpage

\section{Resumo}
Este trabalho apresenta a modelagem e implementacao de um grafo direcionado em C++, com foco em uma entrada padronizada via arquivo texto. A estrutura foi projetada para oferecer operacoes basicas sobre grafos (insercao, remocao e consulta de arestas), leitura robusta de dados e visualizacao da lista de adjacencia. O formato de entrada adotado e:
\begin{verbatim}
n m
u1 v1
u2 v2
...
um vm
\end{verbatim}
em que $n$ representa o numero de vertices, $m$ o numero de arestas e cada par $(u_i, v_i)$ representa uma aresta direcionada $u_i \rightarrow v_i$.

\section{Introducao}
Grafos direcionados sao amplamente utilizados na representacao de dependencias, fluxo de informacao, redes de transporte e problemas de alcance. Em atividades academicas, uma implementacao correta e organizada de grafos e etapa fundamental para evoluir para algoritmos mais avancados (buscas, ordenacao topologica, componentes fortemente conexas, caminhos minimos etc.).

Neste trabalho, o objetivo principal foi estabelecer:
\begin{enumerate}
    \item uma base orientada a objetos para grafos;
    \item um mecanismo padrao de leitura por arquivo;
    \item validacoes de entrada que reduzam erros silenciosos;
    \item uma saida clara em formato de lista de adjacencia.
\end{enumerate}

\section{Objetivos}
\subsection{Objetivo geral}
Implementar uma estrutura de grafo direcionado em C++ com suporte a leitura de instancias por arquivo e impressao em lista de adjacencia.

\subsection{Objetivos especificos}
\begin{enumerate}
    \item Definir e documentar formalmente o formato de entrada.
    \item Implementar carregamento por \texttt{std::istream} e por caminho de arquivo.
    \item Tratar entradas invalidas (cabecalho invalido, arestas fora de faixa, arestas duplicadas, arquivo inacessivel).
    \item Disponibilizar arquivos de teste pequeno, medio e grande.
    \item Disponibilizar \texttt{main.cpp} para validacao pratica.
\end{enumerate}

\section{Especificacao do problema}
\subsection{Formato de entrada}
A entrada segue obrigatoriamente:
\begin{verbatim}
n m
u1 v1
u2 v2
...
um vm
\end{verbatim}

Com as seguintes regras:
\begin{itemize}
    \item $n \ge 0$ e $m \ge 0$;
    \item vertices sao indexados de $0$ a $n-1$;
    \item cada linha posterior contem dois inteiros \texttt{u v};
    \item a aresta e direcionada de \texttt{u} para \texttt{v}.
\end{itemize}

\subsection{Formato de saida}
A impressao do grafo foi padronizada como lista de adjacencia:
\begin{verbatim}
0: 1 2
1: 2
2: 3
3:
\end{verbatim}
Cada linha representa um vertice seguido de seus vizinhos de saida.

\section{Modelagem e arquitetura}
\subsection{Visao geral}
A implementacao foi organizada em tres componentes:
\begin{itemize}
    \item \texttt{Graph}: classe base abstrata com estrutura e operacoes comuns;
    \item \texttt{DirectedGraph}: classe derivada com semantica de grafo direcionado e parser de entrada;
    \item \texttt{main.cpp}: programa de teste para leitura por arquivo e exibicao de resultados.
\end{itemize}

\subsection{Classe \texttt{Graph}}
Responsabilidades:
\begin{itemize}
    \item armazenar lista de adjacencia \texttt{adjacencyList\_};
    \item armazenar contagem de arestas \texttt{edgeCount\_};
    \item fornecer utilitarios de consulta:
    \begin{itemize}
        \item \texttt{vertexCount()}
        \item \texttt{edgeCount()}
        \item \texttt{adjacentTo(vertex)}
        \item \texttt{isValidVertex(vertex)}
        \item \texttt{printAdjacencyList(output)}
    \end{itemize}
\end{itemize}

\subsection{Classe \texttt{DirectedGraph}}
Responsabilidades:
\begin{itemize}
    \item implementar regras de arestas direcionadas:
    \begin{itemize}
        \item \texttt{addEdge(from,to)}
        \item \texttt{removeEdge(from,to)}
        \item \texttt{hasEdge(from,to)}
        \item \texttt{isReachable(source,target)} (DFS para atingibilidade)
    \end{itemize}
    \item implementar carga de dados:
    \begin{itemize}
        \item \texttt{fromStream(std::istream\&)}
        \item \texttt{fromFile(const std::string\&)}
    \end{itemize}
\end{itemize}

\subsection{Programa principal}
O \texttt{main} recebe um argumento obrigatorio (caminho do arquivo) e dois argumentos opcionais (\texttt{source} e \texttt{target}) para consulta de atingibilidade. O programa carrega o grafo e imprime:
\begin{itemize}
    \item confirmacao de leitura;
    \item numero de vertices;
    \item numero de arestas;
    \item lista de adjacencia.
    \item quando \texttt{source} e \texttt{target} sao informados: resultado booleano da pergunta ``Existe caminho de source ate target?''.
\end{itemize}

\section{Detalhes de implementacao}
\subsection{Estrutura de dados}
Foi adotada lista de adjacencia com tipo \texttt{std::vector<std::vector<int>>}. Essa escolha e eficiente para grafos esparsos, reduz consumo de memoria em relacao a matriz de adjacencia e facilita iteracao de vizinhos.

\subsection{Validacoes da leitura}
A leitura implementa as seguintes validacoes:
\begin{enumerate}
    \item falha ao ler \texttt{n} e \texttt{m} gera excecao;
    \item valores negativos para \texttt{n} ou \texttt{m} geram excecao;
    \item falha ao ler qualquer aresta gera excecao;
    \item aresta fora da faixa de vertices gera erro (via \texttt{addEdge});
    \item aresta duplicada gera erro (via \texttt{addEdge});
    \item arquivo inacessivel gera excecao em \texttt{fromFile}.
\end{enumerate}

\subsection{Compatibilidade de codificacao de arquivo}
Foi incluido tratamento para cabecalho BOM UTF-8 no inicio do arquivo, evitando falha na leitura de \texttt{n m} em ambientes como MinGW/Git Bash quando o arquivo foi salvo com BOM.

\subsection{Ordem de insercao e impressao}
A impressao preserva a ordem de insercao dos vizinhos em cada lista. Isso torna a saida deterministica para arquivos de entrada fixos.

\subsection{Verificacao de atingibilidade com DFS}
Foi implementado o metodo \texttt{isReachable(source,target)} na classe \texttt{DirectedGraph}. A estrategia utiliza busca em profundidade (DFS) iterativa com pilha explicita e vetor de visitados.

Fluxo resumido:
\begin{enumerate}
    \item validar se \texttt{source} e \texttt{target} sao vertices validos;
    \item retornar \texttt{true} imediatamente quando \texttt{source == target};
    \item iniciar DFS a partir de \texttt{source};
    \item retornar \texttt{true} assim que \texttt{target} for encontrado;
    \item retornar \texttt{false} se a busca terminar sem encontrar \texttt{target}.
\end{enumerate}

\section{Complexidade}
Seja $n$ o numero de vertices e $m$ o numero de arestas.

\subsection{Operacoes basicas}
\begin{itemize}
    \item \texttt{vertexCount()}: $O(1)$
    \item \texttt{edgeCount()}: $O(1)$
    \item \texttt{adjacentTo(v)}: $O(1)$ para acesso a lista (nao inclui iteracao externa)
\end{itemize}

\subsection{Insercao, remocao e consulta de arestas}
Como cada lista e um \texttt{vector<int>} sem indice auxiliar:
\begin{itemize}
    \item \texttt{addEdge(u,v)}: $O(\deg^+(u))$ (busca para evitar duplicata)
    \item \texttt{removeEdge(u,v)}: $O(\deg^+(u))$
    \item \texttt{hasEdge(u,v)}: $O(\deg^+(u))$
\end{itemize}
No pior caso, essas operacoes podem chegar a $O(n)$ quando um vertice possui alto grau de saida.

\subsection{Consulta de atingibilidade}
Para \texttt{isReachable(source,target)}, no pior caso todos os vertices e arestas alcancaveis a partir de \texttt{source} sao percorridos, resultando em complexidade:
\[
O(n + m)
\]
considerando representacao por lista de adjacencia.

\subsection{Leitura do arquivo}
A leitura processa $m$ arestas. Considerando o custo de verificacao de duplicata em cada insercao, o custo agregado depende da distribuicao de graus de saida. No pior caso extremo, pode chegar a $O(m\cdot n)$. Em cenarios praticos esparsos e com graus moderados, o comportamento e significativamente melhor.

\section{Procedimento experimental}
\subsection{Ambiente}
\begin{itemize}
    \item Linguagem: C++17
    \item Compilador: \texttt{g++}
    \item Execucao validada em terminal Git Bash/PowerShell
\end{itemize}

\subsection{Instancias de teste}
Foram disponibilizadas tres instancias:
\begin{itemize}
    \item \texttt{dados/pequenos/grafo1.txt}
    \item \texttt{dados/medios/grafo1.txt}
    \item \texttt{dados/grandes/grafo1.txt}
\end{itemize}

\subsection{Comandos de compilacao e execucao}
\begin{lstlisting}[style=cppstyle]
g++ -std=c++17 -O2 -Wall -Wextra -o grafo_app.exe src/main.cpp src/graph.cpp src/directed_graph.cpp
\end{lstlisting}

\begin{lstlisting}[style=cppstyle]
./grafo_app.exe dados/pequenos/grafo1.txt
./grafo_app.exe dados/medios/grafo1.txt
./grafo_app.exe dados/grandes/grafo1.txt
\end{lstlisting}

\subsection{Consulta de atingibilidade (evolucao do trabalho)}
\begin{lstlisting}[style=cppstyle]
g++ -std=c++17 -O2 -Wall -Wextra -o grafo_app.exe src/main.cpp src/graph.cpp src/directed_graph.cpp
./grafo_app.exe dados/pequenos/grafo1.txt 0 4
\end{lstlisting}

Saida esperada para o exemplo acima:
\begin{verbatim}
Existe caminho de 0 ate 4? true
\end{verbatim}

\subsection{Resultados esperados}
\begin{itemize}
    \item Pequeno: $n=5$, $m=6$
    \item Medio: $n=12$, $m=20$
    \item Grande: $n=30$, $m=55$
\end{itemize}

Para cada caso, espera-se:
\begin{enumerate}
    \item mensagem de sucesso na leitura;
    \item impressao correta de \texttt{Vertices (n)} e \texttt{Arestas (m)};
    \item lista de adjacencia coerente com o arquivo.
\end{enumerate}

\section{Exemplo de saida}
\begin{verbatim}
Arquivo lido com sucesso: dados/pequenos/grafo1.txt
Vertices (n): 5
Arestas (m): 6
0: 1 2
1: 3 4
2: 3
3: 4
4:
\end{verbatim}

\section{Analise critica}
\subsection{Pontos fortes}
\begin{itemize}
    \item Separacao clara entre estrutura base (\texttt{Graph}) e especializacao (\texttt{DirectedGraph});
    \item Entrada padronizada e documentada;
    \item Leitura por arquivo com tratamento de erro explicito;
    \item Compatibilidade com BOM UTF-8;
    \item verificacao de atingibilidade entre vertices com DFS;
    \item Saida de depuracao/utilizacao simples e objetiva.
\end{itemize}

\subsection{Limitacoes atuais}
\begin{itemize}
    \item custo linear em grau para buscar/remover/consultar arestas;
    \item sem suporte atual a pesos;
    \item ja possui DFS para atingibilidade, mas ainda sem BFS/SCC/ordenacao topologica;
    \item leitura estrita: qualquer inconsistencia interrompe o processamento.
\end{itemize}

\subsection{Melhorias futuras}
\begin{enumerate}
    \item adicionar testes automatizados unitarios;
    \item incluir BFS/DFS e ordenacao topologica;
    \item estudar uso de \texttt{unordered\_set} por vertice para acelerar \texttt{hasEdge}/\texttt{addEdge};
    \item incluir exportacao para formatos padrao (ex.: DOT/Graphviz);
    \item aceitar comentarios em arquivos de entrada.
\end{enumerate}

\section{Conclusao}
O trabalho atingiu os objetivos propostos para esta etapa: formatacao de entrada bem definida, leitura por arquivo implementada, validacoes de consistencia e impressao em lista de adjacencia. A arquitetura escolhida facilita extensao futura para algoritmos classicos de grafos e para analises mais avancadas de desempenho.

Em termos de qualidade de engenharia, o projeto agora possui:
\begin{itemize}
    \item interface de uso clara para execucao;
    \item documentacao de entrada/saida;
    \item conjunto inicial de instancias de teste em diferentes escalas.
\end{itemize}

Isso estabelece uma base solida para evolucao do trabalho nas proximas entregas.

\section*{Checklist para entrega}
\begin{itemize}
    \item Atualizar nome da aluna e dados da disciplina no cabecalho deste arquivo.
    \item Confirmar se o professor exige capa institucional separada.
    \item Compilar e anexar evidencias de execucao (prints ou logs).
    \item Revisar ortografia e padrao exigido pela disciplina.
\end{itemize}

\end{document}
